\section{Problem formulation}
\label{sec:problem}
\providecommand{\AsymptoticsAppendix}{Appendix~\ref{app:asymptotics}}

\subsection{Geometry and initial configuration}
\label{subsec:geometry}

Consider a piecewise homogeneous body under plane strain composed of two
dissimilar isotropic elastic materials joined along a kinked interface
(Figure~\ref{fig:failure-stages}).
Material $i$ has Young's modulus $E_i$ and Poisson's ratio $\nu_i$, where
$i=1,2$ identifies the material. The angle between the two interface rays,
measured through material~2, is $2\alpha$, where $\alpha$ is the half-angle and
\begin{equation*}
  0<\alpha<\pi, \qquad \alpha\ne\frac{\pi}{2}.
\end{equation*}
The value $\alpha=\pi/2$ corresponds to a straight interface and is treated
only as a limiting case.

The origin of the polar coordinate system $(r,\theta)$ is placed at the
interface corner $O$, where $r$ is the distance from $O$ and $\theta$ is
measured from the bisector of the sector occupied by material~2. Symmetry
about this bisector allows the analysis to be restricted to the half-plane
$0\leq\theta\leq\pi$. Material~2 occupies $0<\theta<\alpha$, material~1
occupies $\alpha<\theta<\pi$, and the interface lies along
$\theta=\alpha$. The radial and circumferential displacement components
are denoted by $u_r$ and $u_\theta$, respectively. The circumferential normal
stress is $\sigma_\theta$, and the shear stress in polar coordinates is
$\tau_{r\theta}$. On a radial line, these are the normal and tangential
traction components, respectively.

The loading produces a local field that is symmetric about the bisector.
For the intact interface, the singular part of the stress field along the
interface is
\begin{equation*}
\begin{aligned}
  \tau_{r\theta}(r,\alpha)
    &= Cg_1 r^{\lambda_0}+o\!\left(r^{\lambda_0}\right),\\
  \sigma_\theta(r,\alpha)
    &= Cg_2 r^{\lambda_0}+o\!\left(r^{\lambda_0}\right),
    \qquad r\to0.
\end{aligned}
\end{equation*}
Here, $C$ is the amplitude of the local singular field and incorporates the
effects of the applied loading and the global body geometry; $g_1$ and
$g_2$ are angular coefficients that depend on $\alpha$ and the elastic
properties of the materials; and $\lambda_0\in(-1,0)$ is the singularity
exponent of the intact corner. The notation $o(r^{\lambda_0})$ denotes a
remainder whose ratio to $r^{\lambda_0}$ tends to zero as $r\to0$.
The exponent $\lambda_0$ is the smallest real root of
$D_0(-1-\lambda_0)=0$ in the open interval $(-1,0)$ that is continuously
connected to the physical solution branch. Here, $D_0$ is the characteristic
determinant for the intact corner. Expressions for $D_0$, $g_1$, and $g_2$
are given in \AsymptoticsAppendix.

The initial configuration considered here contains an already existing
V-shaped interfacial shear crack formed by two symmetric branches emanating
from the corner $O$ along the kinked interface. The crack faces are in
frictionless contact. Each branch has length $l$, so the total length of the
two branches is $2l$; below, $l$ is referred to as the crack half-length
(the branch length). The contact state is consistent with
\begin{equation*}
  Cg_2<0,
\end{equation*}
for which the normal stress along the interface is compressive. The
V-shaped interfacial crack is treated as a prescribed stationary
configuration and $l$ as a prescribed fixed parameter. The applied loading
is assumed insufficient to cause loss of local stability at the crack-branch
tips. Interfacial propagation from those tips is outside the local problem
near $O$ and is considered below only as a competing fracture mechanism.

The prescribed interfacial crack retains a power-law stress singularity at
the corner~\cite{NazarenkoKipnis2022Stress}. This field can support a small
tensile process zone in one of the materials; its characteristics are used
below to formulate the conditions for branching initiation at the corner. Under the assumed geometric and loading symmetry, the bisector of the sector
occupied by that material is a natural direction for the zone. Along this
direction, the local Mode~II (in-plane shear) component at the zone tip
vanishes by symmetry, consistent with the principle of local
symmetry~\cite{BarenblattCherepanov1961,GoldsteinSalganik1974}. Zones with
arbitrary orientations are not considered. The bisector direction is
therefore part of the assumed symmetric model and is not obtained by
separately optimizing the branching direction.

Let $\beta_i$ denote the polar angle of the bisector in material~$i$:
\begin{equation*}
  \beta_1=\pi,\qquad \beta_2=0.
\end{equation*}
The subscript $i$ in $\beta_i$, $d_i$, $\sigma_i$, and $Q_i$ identifies
the material throughout.

The process zone in material~$i$ is modeled as a normal displacement
discontinuity of length $d_i$ subjected to a constant normal cohesive
traction $\sigma_i>0$, which represents the material's resistance to
separation. Figure~\ref{fig:failure-stages} shows the zone in material~1
as a solid line and the alternative zone in material~2 as a dashed line.
These are alternative configurations, considered separately. For each
admissible direction, the solution defines a quasi-static family of
equilibrium process-zone states; the zone opening and energy characteristic
are then used to formulate branching-initiation criteria.

\begin{figure}[htbp]
  \centering
  % Vector redraw of the author's original Figure 1.
% The figure is intentionally kept as TikZ so that all labels are typeset
% by LaTeX and remain vector objects in the compiled manuscript PDF.
\begin{tikzpicture}[x=1cm,y=1cm,>=Stealth,line cap=round,line join=round]
\tikzset{
  panel/.style={draw=black,line width=.6pt,rounded corners=6pt},
  interface/.style={draw=black,line width=.7pt},
  axis/.style={draw=black,line width=.45pt},
  crack/.style={draw=black,double=white,double distance=2.0pt,line width=.45pt},
  processzone/.style={draw=black,line width=2.0pt},
  alternativezone/.style={draw=black,line width=1.2pt,dashed}
}

% Common panel geometry.  The interface rays are symmetric about the
% horizontal bisector of material 2, as in the author's original sketch.
\newcommand{\failurepanel}[2]{%
  \begin{scope}[xshift=#1cm]
    \draw[panel] (-2.20,-1.55) rectangle (2.20,1.55);
    \coordinate (O) at (0,0);
    \coordinate (Iu) at ($(O)+(118:2.05)$);
    \coordinate (Il) at ($(O)+(242:2.05)$);
    \draw[interface] (O)--(Iu);
    \draw[interface] (O)--(Il);
    \draw[axis] (-2.05,0)--(2.05,0);

    % Polar coordinates and the half-angle alpha.
    \draw[->,line width=.65pt] (O)--++(27:1.62)
      node[above right=-1pt,font=\small] {$r$};
    \draw[->,line width=.5pt] (0.58,0)
      arc[start angle=0,end angle=27,radius=.58];
    \node[font=\small,fill=white,inner sep=.7pt] at (0.93,0.18) {$\theta$};
    \draw[->,line width=.5pt] (0.83,0)
      arc[start angle=0,end angle=118,radius=.83];
    \node[font=\small,fill=white,inner sep=.7pt] at (0.16,0.88) {$\alpha$};

    \node[font=\small,fill=white,inner sep=.7pt] at (0.18,-0.20) {$O$};
    \node[font=\small] at (-1.45,-0.92) {$E_1,\,\nu_1$};
    \node[font=\small] at ( 1.30,-0.92) {$E_2,\,\nu_2$};
    #2
  \end{scope}%
}

% (a) Intact broken interface.
\failurepanel{0}{%
  \node[font=\small] at (0,-1.88) {(a)};
}

% (b) Symmetric interfacial shear cracks of length l.
\failurepanel{5.00}{%
  \draw[crack] (O)--++(118:1.30);
  \draw[crack] (O)--++(242:1.30);
  \node[font=\small] at (-0.77, 1.18) {$l$};
  \node[font=\small] at (-0.77,-1.18) {$l$};
  \node[font=\small] at (0,-1.88) {(b)};
}

% (c) Mode-I process-zone alternatives along the two material bisectors.
\failurepanel{10.00}{%
  \draw[crack] (O)--++(118:1.30);
  \draw[crack] (O)--++(242:1.30);
  \node[font=\small] at (-0.77, 1.18) {$l$};
  \node[font=\small] at (-0.77,-1.18) {$l$};
  \draw[processzone] (O)--(-1.35,0);
  \draw[alternativezone] (O)--(1.35,0);
  \node[font=\small,below=2pt] at (-0.75,0) {$d_1$};
  \node[font=\small,below=2pt] at ( 0.75,0) {$d_2$};
  \node[font=\small] at (0,-1.88) {(c)};
}
\end{tikzpicture}

  \caption{Local configurations near the interface corner:
  (a) intact kinked interface; (b) prescribed stationary V-shaped
  interfacial shear crack with two branches of length $l$ each (total branch
  length $2l$); (c) process zone of length $d_1$
  in material~1 and the alternative zone of length $d_2$ in material~2
  (dashed line). The panels compare the corresponding local configurations
  and do not prescribe the order in which they occur.}
  \label{fig:failure-stages}
\end{figure}

\subsection{Local problem and boundary conditions}
\label{subsec:boundary-conditions}

Let $L$ be a characteristic outer dimension of the body. We assume the
separation of scales
\begin{equation*}
  d_i\ll l\ll L
\end{equation*}
and neglect friction between the interfacial crack faces. Within this
quasi-static formulation, the outer-field amplitude $C$ may vary, whereas
the crack half-length $l$, and hence the total length $2l$ of the two
branches, remain fixed. Changes in loading alter the equilibrium zone length
$d_i$ and the associated local characteristics. In the local problem, the
piecewise homogeneous body is replaced by an infinite plane, and the two
branches of the V-shaped interfacial crack are treated as semi-infinite.
Their finite length and the applied loading enter through an asymptotic
matching condition~\cite{Nayfeh1973Perturbation}. This is a one-way local
approximation: the known field of the stationary interfacial crack loads the
inner problem near $O$, whereas the feedback of the small process zone on
the fields at the crack-branch tips is not determined.

For a field quantity $f$ defined on both sides of the ray
$\theta=\alpha$, its jump across the interface is defined by
\begin{equation*}
  \jump{f}
  =f(r,\alpha+0)-f(r,\alpha-0),
\end{equation*}
where $\alpha+0$ and $\alpha-0$ denote the limiting values approached
from materials~1 and~2, respectively. Symmetry, traction continuity, and
frictionless contact give
\begin{equation}
\label{eq:interface-boundary-conditions}
\begin{aligned}
  &\theta\in\{0,\pi\}: &&\tau_{r\theta}=0,\\
  &\theta=\alpha: &&
    \jump{\sigma_\theta}=0,
    \quad \jump{\tau_{r\theta}}=0,
    \quad \tau_{r\theta}=0,
    \quad \jump{u_\theta}=0.
\end{aligned}
\end{equation}
The condition $\jump{u_\theta}=0$ imposes normal contact between the crack
faces, whereas $\tau_{r\theta}=0$ expresses the absence of friction.

On the bisector of the material sector containing the process zone, the
following mixed boundary conditions apply:
\begin{equation}
\label{eq:process-zone-boundary-conditions}
\begin{aligned}
  &\theta=\beta_i,\quad 0<r<d_i:
    &&\sigma_\theta(r,\beta_i)=\sigma_i,\\
  &\theta=\beta_i,\quad r>d_i:
    &&u_\theta(r,\beta_i)=0,\\
  &\theta=\beta_j,\quad r>0,\quad j\ne i:
    &&u_\theta(r,\beta_j)=0.
\end{aligned}
\end{equation}
Here, $j=3-i$ identifies the other material. Thus, the symmetry
condition on $u_\theta$ is replaced by a prescribed constant normal
traction only over $0<r<d_i$ along the active direction
$\theta=\beta_i$.

At intermediate distances $d_i\ll r\ll l$, the solution must match the
asymptotic field of the interfacial crack in the absence of a process
zone~\cite{NazarenkoKipnis2022Stress}. In the inner problem, this matching
condition takes the form
\begin{equation}
\label{eq:matched-far-field}
  \sigma_\theta(r,\beta_i)
    =CQ_i r^\lambda+o\!\left(r^{-1}\right),
    \qquad r\to\infty.
\end{equation}
Here, $Q_i$ is the amplitude coefficient for material~$i$, which depends
on the interface angle, the elastic properties, and the branch length
$l$, as specified in \AsymptoticsAppendix. The exponent
$\lambda\in(-1,0)$ characterizes the corner singularity in the
configuration with the prescribed stationary V-shaped interfacial crack. It is the smallest real root of
$D(-1-\lambda)=0$ in $(-1,0)$ that is continuously connected to the
physical branch, where $D$ is the characteristic determinant for the
cracked configuration. Expressions for $D$ and $Q_i$ are given in
\AsymptoticsAppendix. The limit $r\to\infty$ is understood in the
inner-scale sense of asymptotic matching, rather than as the far field
of the finite body; the remainder $o(r^{-1})$ decays faster than $r^{-1}$
in that limit.

For the subsequent Wiener--Hopf construction, we use superposition as in
the related problem~\cite{KaminskyDudykPolishchuk2024}. We subtract the known field of the stationary V-shaped interfacial crack
without a process zone and introduce the correction normal stress
\begin{equation*}
  \widehat{\sigma}_i(r)
  =\sigma_\theta(r,\beta_i)-CQ_i r^\lambda.
\end{equation*}
The matching condition then gives
\begin{equation*}
  \widehat{\sigma}_i(r)=o(r^{-1}),
  \qquad r\to\infty,
\end{equation*}
while on the process-zone segment $0<r<d_i$,
\begin{equation*}
  \widehat{\sigma}_i(r)
  =\sigma_i-CQ_i r^\lambda.
\end{equation*}
The field of the stationary V-shaped crack without a process zone satisfies
the bisector symmetry condition $u_\theta=0$, so subtracting it does not alter the displacement-gradient
transform used below. The functional problem is therefore formulated for the correction field,
with the known power-law term included in the prescribed part.

A process zone can form in material~$i$ only if the normal stress on its
bisector is tensile. Equation~\eqref{eq:matched-far-field} therefore gives
the physical admissibility condition
\begin{equation*}
  CQ_i>0.
\end{equation*}
The model thus requires both $Cg_2<0$, for contact between the
interfacial crack faces, and $CQ_i>0$, for formation of a process zone
in material~$i$. For the parameter sets examined below, the numerical
analysis shows that these conditions can hold simultaneously for only one
of the two materials; only that material admits a physically admissible
process zone.

\subsection{Asymptotic fields near the process-zone tip}
\label{subsec:zone-tip-asymptotics}

Near the outer tip of the displacement discontinuity, the stress and normal
displacement gradient have the standard inverse-square-root asymptotics for a
homogeneous material:
\begin{equation}
\label{eq:zone-tip-asymptotics}
\begin{aligned}
  &\theta=\beta_i,\quad r\to d_i+0:
  &&\sigma_\theta(r,\beta_i)
    \sim \frac{k_i}{\sqrt{2\pi(r-d_i)}},\\[3pt]
  &\theta=\beta_i,\quad r\to d_i-0:
  &&\frac{\partial u_\theta}{\partial r}(r,\beta_i)
    \sim -\frac{2(1-\nu_i^2)}{E_i}
    \frac{k_i}{\sqrt{2\pi(d_i-r)}}.
\end{aligned}
\end{equation}
Here, $k_i$ is the local Mode~I stress intensity factor at the process-zone
tip; $r\to d_i+0$ and $r\to d_i-0$ denote approach from ahead of the
zone and from within it, respectively. The boundary-value problem is closed by
imposing the standard stress-regularity condition for this
model~\cite{KaminskyDudykPolishchuk2024},
\begin{equation*}
  k_i=0.
\end{equation*}
This condition removes the inverse-square-root singular term and is
used below to determine the equilibrium zone length $d_i$. It does not
require the bounded part of the stress at the zone tip to vanish.
